\documentclass{article}
\usepackage[margin=1in]{geometry}
\usepackage{longtable}
\begin{document}

\section*{Phase 10 Plan - Single Owner}

\textbf{Source of truth:} \texttt{Documentation/SRS.tex}, especially Section 2.5 ``Implementation Readiness Checklist'', Section 2.6 ``Definition of v1 Complete'', Section 7 ``Evidence, Alerting, and Analyst Workflow'', Section 8 ``Replay, Backtesting, and Validation'', Section 9 ``Operations, Security, and Governance'', Section 10 ``Cost Model and Budget'', and Section 11 ``Roadmap and Stage Gates'', plus the Phase 8 and Phase 9 closeout artifacts under \texttt{Documentation/phases/} and \texttt{reports/phase9/}.\\
\textbf{Execution status:} Planned single-owner follow-on after the Phase 9 closeout refresh.\\
\textbf{Entry assumption:} Phase 9 already produced a materially populated local evidence packet through Phase 6, but the repo still does \textbf{not} qualify as SRS-complete v1 because the current evidence path is still noop-provider-backed and the current boosted-tree model evidence is still too small and train-only to support a held-out-strength claim.

\textbf{Phase goal:} Convert the repo from \textbf{materially populated but not SRS-complete v1} into a \textbf{defensible, operator-hardened, real-provider-validated v1 package} that can satisfy the SRS completion contract without caveating away evidence quality.\\
\textbf{Interpretation note:} The formal SRS roadmap ends at Phase 7. Phase 8 handled consolidation and truthfulness, and Phase 9 handled evidence materialization. Phase 10 is therefore a \textbf{single-owner operational-hardening and evidence-quality phase} derived from the exact blockers left open by Phase 9 and from the mandatory SRS requirements that remain stronger than the current committed proof.\\
\textbf{Implementation note:} Phase 10 must prefer stronger evidence quality, held-out validation, operator readiness, and governance enforcement over adding fresh research breadth.

\section*{Why Phase 10 Exists}
Phase 9 solved the ``empty artifacts'' problem but did not yet solve the ``credible final claim'' problem.
\begin{itemize}
\item The repo now has a replay-linked local packet from candidate generation through alerting, validation, and shadow ML.
\item The current Phase 4 evidence path is still based on local seeded or noop-provider behavior rather than a stronger real-provider-backed evidence packet.
\item The current Phase 6 LightGBM artifact is still tiny and train-only, so it cannot honestly carry the SRS requirement for meaningful model evaluation against wallet-unaware baselines.
\item The SRS still requires monitoring, secrets discipline, redaction policy, analyst-feedback capture, and budget guardrails to be treated as first-class runtime requirements.
\item The remaining work is therefore not another speculative subsystem. It is the final conversion from a believable prototype package into a defensible operational v1.
\end{itemize}

\section*{Owned Deliverables (Synthesized From SRS + Phase 9 Findings)}
\begin{enumerate}
\item Replace the canonical Phase 4 noop-provider path with one real-provider-backed evidence workflow that records timestamped queries, results, cache behavior, and budget usage for replay-linked alerts.
\item Materialize a stronger analyst-feedback corpus across more than one alert window so suppression, severity policy, and usefulness labels are no longer supported by only a single example path.
\item Produce Phase 5 validation on a held-out-strength window family with exact time-split or category-holdout boundaries that match the SRS validation contract.
\item Produce Phase 6 boosted-tree evaluation on a meaningfully larger dataset with required wallet-unaware baselines, held-out reporting, calibration outputs, and threshold recommendations that are not train-only.
\item Add the mandatory SRS operational controls: monitoring coverage, archive and database integrity checks, incident runbooks, backup discipline, and budget guardrails for evidence-provider usage.
\item Add the mandatory SRS security and governance controls: secret storage policy, role-scoped DB usage where applicable, wallet redaction in user-facing alerts, and deployment or model-version audit logging.
\item Write one promotion memo stating whether the canonical operating mode remains \texttt{rule\_based\_plus\_shadow\_ml} or whether any part of the model stack is now strong enough for a stricter operational role.
\item Produce one final completion memo that answers the SRS v1-complete question with the strongest honest answer the evidence supports.
\end{enumerate}

\section*{Phase 10 Exit Criteria}
\begin{itemize}
\item At least one canonical alert-evidence packet is backed by real providers rather than noop-provider placeholders.
\item Analyst feedback and suppression behavior are supported by repeated examples, not only a single seeded demonstration.
\item Phase 5 validation includes held-out-strength reporting that can be cited without leakage caveats.
\item Phase 6 evaluation includes a boosted-tree ranker compared against wallet-unaware baselines on non-trivial held-out data with saved calibration outputs.
\item Mandatory SRS monitoring, incident, security, governance, and budget guardrails are documented and materially implemented for the canonical v1 path.
\item The repo can answer either ``SRS-complete v1'' or ``still blocked'' with a small, explicit blocker list rather than broad uncertainty.
\end{itemize}

\section*{Locked Scope Decisions}
\begin{itemize}
\item Phase 10 is a hardening and completion phase, not a new Phase 7 research branch.
\item Canonical v1 remains local-first unless a separately documented architecture decision proves a stronger path.
\item No autonomous execution, live capital deployment, or cloud-first redesign belongs in this phase.
\item A real-provider-backed evidence packet must respect explicit daily and monthly provider budgets.
\item Train-only wins do not count as final model proof.
\item If held-out evidence is still weak after Phase 10, the repo must say so plainly rather than promoting weaker claims.
\item Governance language from the SRS remains binding: no user-facing output may call a wallet or cluster an insider.
\end{itemize}

\section*{Concrete Artifacts Required}
\begin{longtable}{p{0.24\linewidth}p{0.36\linewidth}p{0.26\linewidth}}
\textbf{Workstream} & \textbf{Artifact} & \textbf{Done Condition} \\
Real-provider evidence & Provider-backed alert packet and evidence review memo & Queries, snapshots, cache metadata, budget usage, and alert payload all exist for a canonical window \\
Analyst loop expansion & Multi-alert feedback summary and suppression review & Multiple alerts have persisted analyst actions and a documented false-positive or usefulness review path \\
Held-out validation & Time-split or category-holdout validation pack & Phase 5 reporting uses exact held-out boundaries and can be defended against leakage concerns \\
Phase 6 completion & Larger held-out LightGBM or CatBoost evaluation bundle & Baselines, calibration, thresholds, and model card exist on non-trivial held-out data \\
Ops hardening & Monitoring and incident runbook bundle & SRS-required metrics, integrity checks, and failure-response workflows are documented and wired into the canonical path \\
Security and governance & Redaction, secrets, audit, and policy memo & User-facing alerts redact wallet identifiers by default and operator-sensitive actions are governed explicitly \\
Final decision & Phase 10 completion memo & The repo states whether v1 is complete and what exact blocker remains if not \\
\end{longtable}

\section*{Canonical Phase 10 Workflow}
\begin{enumerate}
\item Choose the canonical held-out window family and the canonical live or replay-backed alert windows that will support the final claim.
\item Replace or upgrade the current evidence-provider path so the chosen alert packets use real provider calls, timestamps, caching, and budget accounting.
\item Run the alert loop repeatedly enough to evaluate suppression, severity, and analyst usefulness on more than one example.
\item Run held-out Phase 5 validation on the same feature and label contracts used by the canonical Phase 6 evaluation.
\item Train and evaluate the boosted-tree ranker on a larger dataset with held-out reporting, calibration, and baseline comparisons.
\item Implement or verify the mandatory SRS operations, security, governance, and budget controls on the canonical path.
\item Write the operating-mode promotion memo and the final completion memo.
\end{enumerate}

\section*{Concrete Execution Sequence}
\textbf{Block 1: Evidence-provider hardening}
\begin{itemize}
\item Replace noop-provider defaults on the canonical Phase 4 path with real-provider-backed retrieval where feasible.
\item Log provider, query time, timeout, cache hit state, and spend or quota metadata for each evidence fetch.
\item Produce one review packet that shows the evidence layer is no longer only a seeded local demonstration.
\end{itemize}

\textbf{Block 2: Analyst-loop expansion and suppression tuning}
\begin{itemize}
\item Materialize more than one alert episode through the canonical delivery path.
\item Capture structured analyst actions, follow-up notes, and resolution outcomes where available.
\item Review duplicate suppression, severity inflation, and false-positive handling against real alert traffic.
\end{itemize}

\textbf{Block 3: Held-out replay and model evaluation}
\begin{itemize}
\item Freeze exact held-out boundaries before rerunning Phase 5 and Phase 6 evaluation.
\item Regenerate replay-derived training or evaluation inputs without future leakage.
\item Compare the boosted-tree ranker against the required wallet-unaware baselines on non-trivial held-out data.
\item Update calibration, threshold recommendations, and the model card from held-out outputs rather than train-only summaries.
\end{itemize}

\textbf{Block 4: Ops, security, and governance closeout}
\begin{itemize}
\item Verify monitoring coverage for collector health, archive writes, detector backlog, Redis usage, PostgreSQL latency, candidate rate, alert rate, and evidence latency.
\item Consolidate incident runbooks for collector death, reconnect storms, disk pressure, schema drift, and daily integrity checks.
\item Enforce secret handling, wallet redaction, deployment logging, and evidence budget caps on the canonical path.
\end{itemize}

\textbf{Block 5: Final completion decision}
\begin{itemize}
\item Decide whether the canonical operating mode stays \texttt{rule\_based\_plus\_shadow\_ml} or whether stronger promotion is justified.
\item Refresh the closeout documents with the new evidence-quality and hardening results.
\item Write the final Phase 10 memo with either a complete SRS-v1 claim or a minimal explicit blocker list.
\end{itemize}

\section*{Expected Areas To Touch}
\begin{itemize}
\item \texttt{Documentation/phases/phase10.tex}
\item \texttt{Documentation/INDEX.tex} and \texttt{Documentation/phases/README.tex}
\item \texttt{phase4/} evidence-provider integrations and alert rendering paths
\item \texttt{phase5/}, \texttt{phase6/}, and their runner entrypoints for held-out evaluation
\item \texttt{database/*RUNBOOK.md} and related operator documentation
\item generated outputs under \texttt{reports/phase10/}, \texttt{reports/phase5/}, \texttt{reports/phase6/}, and refreshed closeout artifacts
\end{itemize}

\section*{Non-Goals}
\begin{itemize}
\item No new graph-model or deep-sequence research branch as part of this phase.
\item No broad architecture rewrite that resets the proven local-first path.
\item No attempt to inflate evidence quality by widening the claim while keeping tiny datasets.
\item No hiding of provider limitations, quota failures, or weak alert quality behind more polished docs.
\end{itemize}

\section*{Phase 10 Readiness Standard}
Phase 10 is complete when the repo can defend a real-provider-backed alert path, a held-out-strength validation path, and the mandatory SRS operations and governance controls as one coherent local-first v1 system. At that point the project is no longer only materially populated. It is either honestly SRS-complete v1 or close enough that the remaining blocker list is narrow, concrete, and operational rather than structural.

\end{document}
